\chapter{Future Work}
\label{ch:future}

\begin{itemize}
\item \textbf{Further algorithms.} The \ac{GA} of this milestone plugs into the pipeline unchanged;
the next milestone adds a swarm method (\ac{PSO}) the same way, compared on the same case studies with
the same budget, seeds and statistical tests.
\item \textbf{Hybrid and parallel \ac{GA}.} A final local search with \ac{SA}'s move around the best
design would add \ac{SA}'s fine-tuning to the \ac{GA}'s exploration, and evaluating the children of a
generation in parallel would make the \ac{GA} several times faster in wall-clock time.
\item \textbf{Trade-off without weights.} The weights decide how much one more spacecraft or one more
day is worth. \ac{NSGA}-II will return the whole Pareto front of coverage, fleet cost, energy and
mission time, on which the weighted designs of this milestone should lie.
\item \textbf{Gravity.} Orbits are ideal Kepler circles. The irregular gravity of Eros, and even more
of the dog-bone shaped Kleopatra, makes close orbits drift or become unstable; a gravity field
computed from the shape (polyhedron model) is the next step towards realism.
\item \textbf{Seasons and pointing.} The Sun direction is fixed and the cameras always point at the
centre of mass; a moving Sun and off-nadir pointing would add realism and decision variables.
\item \textbf{Statistics.} Some forced-size cases use three runs; more runs would tighten the spreads
and the $p$-values.
\item \textbf{Reproducible validation.} Every change to the objective, constraints or mesh should
trigger a regeneration of the comparison matrix with the same seeds and locked dependencies. The
validation output should retain the git commit, package versions, mesh resolution, objective weights
and reference fleet size alongside each result.
\item \textbf{Sensitivity of the trade-off.} The fixed-reference normalization correctly penalises
fleets larger than ten, but the preferred fleet still depends on the weights. A systematic sensitivity
study, followed by a Pareto-front analysis, would quantify how robust the selected mission is to the
relative value assigned to coverage, cost, energy and duration.
\end{itemize}
